\chapter{Source Code Appendix}
\section{main.py}
\begin{lstlisting}[language=Python]
from fastapi import FastAPI, File, Query, UploadFile, HTTPException, Form
from fastapi.responses import FileResponse, PlainTextResponse
import uvicorn
import joblib
import numpy as np
from pydantic import BaseModel
import pandas as pd
from datetime import datetime


app = FastAPI(
    title="Credit Card Fraud Detection API",
    description="""An API that utilizes a Machine Learning model that detects if a credit card transaction is fraudulent or not based on the following features: hours, amount, transaction type etc.""",
    version="1.0.0", debug=True)

# --- Runtime Counters ---
_startup_time: datetime = datetime.utcnow()
_stats = {
    "total_predictions": 0,
    "fraud_count": 0,
    "safe_count": 0,
    "total_probability": 0.0,
}


# Load Model and Artifacts
try:
    import os
    # Get the directory of the current script
    current_dir = os.path.dirname(os.path.abspath(__file__))
    # Construct absolute path to the model
    model_path = os.path.join(current_dir, '..', 'models', 'credit_fraud.pkl')

    artifacts = joblib.load(model_path)
    # Check if it's the new dictionary format or old model
    if isinstance(artifacts, dict) and 'model' in artifacts:
        model = artifacts['model']
        THRESHOLD = artifacts.get('threshold', 0.5)
        ISO_FOREST = artifacts.get('iso_forest', None)
        ISO_FEATURES = artifacts.get('iso_features', [])
        TYPE_ENCODING = artifacts.get('type_encoding', {})
        GLOBAL_FRAUD_RATE = artifacts.get('global_fraud_rate', 0.0)
        MODEL_FEATURES = artifacts.get('features', [])
        MODEL_VERSION = artifacts.get('model_version', 'v2.0')
        print(f"Loaded {MODEL_VERSION} model with threshold: {THRESHOLD}")
    else:
        model = artifacts
        THRESHOLD = 0.5
        ISO_FOREST = None
        ISO_FEATURES = []
        TYPE_ENCODING = {}
        GLOBAL_FRAUD_RATE = 0.0
        MODEL_FEATURES = []
        MODEL_VERSION = 'legacy'
        print("Loaded legacy model (threshold 0.5)")
except Exception as e:
    print(f"Error loading model: {e}")
    model = None
    THRESHOLD = 0.5
    ISO_FOREST = None
    ISO_FEATURES = []
    TYPE_ENCODING = {}
    GLOBAL_FRAUD_RATE = 0.0
    MODEL_FEATURES = []
    MODEL_VERSION = 'none'

@app.get("/", response_class=PlainTextResponse)
async def running():
  note = """
Credit Card Fraud Detection API 🙌🏻

Model: Advanced Stacking Ensemble (RF + XGB + LightGBM + CatBoost → LR meta)
Note: add "/docs" to the URL to get the Swagger UI Docs or "/redoc"
  """
  return note

favicon_path = 'favicon.png'
@app.get('/favicon.png', include_in_schema=False)
async def favicon():
    return FileResponse(favicon_path)
																	
class fraudDetection(BaseModel):
    step:int
    types:int
    amount:float	
    oldbalanceorig:float	
    newbalanceorig:float	
    oldbalancedest:float	
    newbalancedest:float	
    isflaggedfraud:float


def _build_features(data: "fraudDetection") -> pd.DataFrame:
    """
    Central feature engineering function — mirrors train.py exactly.
    Returns a DataFrame with the correct column order for the loaded model.
    """
    # ── Temporal ─────────────────────────────────────────────────────────────
    hour        = data.step % 24
    day_of_week = (data.step // 24) % 7
    is_night    = 1 if hour <= 6 else 0

    # ── Balance errors ────────────────────────────────────────────────────────
    errorBalanceOrig = data.newbalanceorig + data.amount - data.oldbalanceorig
    errorBalanceDest = data.oldbalancedest + data.amount - data.newbalancedest

    # ── Log-amount ────────────────────────────────────────────────────────────
    amount_log = np.log1p(data.amount)

    # ── Advanced ratio features ───────────────────────────────────────────────
    amount_to_balance_ratio    = data.amount / (data.oldbalanceorig + 1)
    balance_drain_flag         = 1 if data.newbalanceorig == 0 else 0
    dest_balance_increase_ratio = (
        (data.newbalancedest - data.oldbalancedest) / (data.amount + 1)
    )
    is_round_amount  = 1 if data.amount % 1000 == 0 else 0
    balance_dest_ratio = data.amount / (data.oldbalancedest + 1)

    # ── Target-encoded 'type' ─────────────────────────────────────────────────
    type_encoded = TYPE_ENCODING.get(data.types, GLOBAL_FRAUD_RATE)

    # ── Isolation Forest anomaly score ────────────────────────────────────────
    isolation_score = 0.0
    if ISO_FOREST is not None and ISO_FEATURES:
        iso_input = pd.DataFrame([{
            "amount":                    data.amount,
            "amount_log":               amount_log,
            "amount_to_balance_ratio":  amount_to_balance_ratio,
            "errorBalanceOrig":         errorBalanceOrig,
            "errorBalanceDest":         errorBalanceDest,
            "balance_drain_flag":       balance_drain_flag,
            "dest_balance_increase_ratio": dest_balance_increase_ratio,
        }])[ISO_FEATURES]
        isolation_score = float(-ISO_FOREST.decision_function(iso_input)[0])

    # ── Assemble feature dict in the EXACT column order from training ─────────
    # This order must match FINAL_FEATURES in train.py
    features_dict = {
        "step":                        [data.step],
        "type":                        [data.types],  # kept as numeric fallback
        "amount":                      [data.amount],
        "oldbalanceOrg":               [data.oldbalanceorig],
        "newbalanceOrig":              [data.newbalanceorig],
        "oldbalanceDest":              [data.oldbalancedest],
        "newbalanceDest":              [data.newbalancedest],
        "isFlaggedFraud":              [data.isflaggedfraud],
        "hour":                        [hour],
        "day_of_week":                 [day_of_week],
        "is_night":                    [is_night],
        "errorBalanceOrig":            [errorBalanceOrig],
        "errorBalanceDest":            [errorBalanceDest],
        "amount_log":                  [amount_log],
        "amount_to_balance_ratio":     [amount_to_balance_ratio],
        "balance_drain_flag":          [balance_drain_flag],
        "dest_balance_increase_ratio": [dest_balance_increase_ratio],
        "is_round_amount":             [is_round_amount],
        "balance_dest_ratio":          [balance_dest_ratio],
        "type_encoded":                [type_encoded],
        "isolation_score":             [isolation_score],
    }

    df = pd.DataFrame(features_dict)

    # If the model has a specific feature list, reorder/subset to match it
    if MODEL_FEATURES:
        # Keep only columns the model was trained on, in the right order
        available = [c for c in MODEL_FEATURES if c in df.columns]
        df = df[available]

    return df


@app.post('/predict')
def predict(data: fraudDetection):
    if model is None:
        return {"error": "Model not loaded"}

    features = _build_features(data)

    # Predict
    probability = model.predict_proba(features)[0][1]

    # Apply optimised threshold
    prediction = 1 if probability >= THRESHOLD else 0

    # Update counters
    _stats["total_predictions"] += 1
    _stats["total_probability"] += float(probability)
    if prediction == 1:
        _stats["fraud_count"] += 1
        return {"result": "Fraudulent", "probability": float(probability)}
    else:
        _stats["safe_count"] += 1
        return {"result": "Not Fraudulent", "probability": float(probability)}

# -----------------
# Health & Stats Endpoints
# -----------------

@app.get("/health")
async def health():
    """Returns API health, uptime, and model status."""
    uptime_seconds = (datetime.utcnow() - _startup_time).total_seconds()
    hours, remainder = divmod(int(uptime_seconds), 3600)
    minutes, seconds = divmod(remainder, 60)
    return {
        "status": "online",
        "model_loaded": model is not None,
        "threshold": float(THRESHOLD),
        "uptime": f"{hours:02d}h {minutes:02d}m {seconds:02d}s",
        "uptime_seconds": uptime_seconds,
    }


@app.get("/stats")
async def stats():
    """Returns running prediction statistics."""
    total = _stats["total_predictions"]
    avg_prob = (_stats["total_probability"] / total) if total > 0 else 0.0
    fraud_rate = (_stats["fraud_count"] / total * 100) if total > 0 else 0.0
    return {
        "total_predictions": total,
        "fraud_count": _stats["fraud_count"],
        "safe_count": _stats["safe_count"],
        "avg_probability": round(avg_prob, 4),
        "fraud_rate_pct": round(fraud_rate, 2),
    }


# -----------------
# Explainable AI (SHAP)
# -----------------
import shap

# Initialize Explainer (Lazy Load to prevent startup crash)
explainer = None

def get_explainer():
    global explainer
    if explainer is not None:
        return explainer

    try:
        if model:
            # v3.0 stacking: try LightGBM first (fastest SHAP), then XGBoost, then RF
            for estimator_key in ['lgbm', 'xgb', 'rf']:
                try:
                    est = model.named_estimators_.get(estimator_key)
                    if est is None:
                        continue
                    explainer = shap.TreeExplainer(est)
                    print(f"SHAP Explainer initialised using '{estimator_key}'.")
                    return explainer
                except Exception:
                    continue

            # Legacy fallback (v1/v2 VotingClassifier)
            xgb_est = model.named_estimators_.get('xgb')
            if xgb_est is not None:
                explainer = shap.TreeExplainer(xgb_est)
                print("SHAP Explainer (legacy): using XGBoost from VotingClassifier.")
                return explainer

            raise ValueError("No supported estimator found for SHAP.")
    except Exception as e:
        print(f"Error initialising SHAP explainer: {e}")
        return None


@app.post('/explain')
def explain(data: fraudDetection):
    explainer_instance = get_explainer()
    if explainer_instance is None:
        return {"error": "Explainer could not be initialized"}

    features = _build_features(data)

    # Calculate SHAP values
    try:
        shap_values = explainer_instance.shap_values(features)
        # For binary classification some explainers return a list
        if isinstance(shap_values, list):
            shap_arr = shap_values[1][0]  # class 1
        else:
            shap_arr = shap_values[0]

        return {
            "base_value":    float(explainer_instance.expected_value
                                  if not isinstance(explainer_instance.expected_value, (list, np.ndarray))
                                  else explainer_instance.expected_value[1]),
            "shap_values":   shap_arr.tolist(),
            "feature_names": list(features.columns),
        }
    except Exception as e:
        print(f"SHAP Params Error: {e}")
        return {"error": f"Calculation failed: {str(e)}"}


if __name__ == "__main__":
    uvicorn.run(app, host="0.0.0.0", port=8000)

\end{lstlisting}

\section{app.py}
\begin{lstlisting}[language=Python]
from fastapi import FastAPI, File, Query, UploadFile, HTTPException, Form
from fastapi.responses import FileResponse, PlainTextResponse
import uvicorn
import joblib
import numpy as np
from pydantic import BaseModel



app = FastAPI(
    title="Credit Card Fraud Detection API",
    description="""An API that utilises a Machine Learning model that detects if a credit card transaction is fraudulent or not based on the following features: hours, amount, transaction type etc.""",
    version="1.0.0", debug=True)


model = joblib.load('credit_fraud.pkl')

@app.get("/", response_class=PlainTextResponse)
async def running():
  note = """
Credit Card Fraud Detection API 🙌🏻

Note: add "/docs" to the URL to get the Swagger UI Docs or "/redoc"
  """
  return note

favicon_path = 'favicon.png'
@app.get('/favicon.png', include_in_schema=False)
async def favicon():
    return FileResponse(favicon_path)
																	
class fraudDetection(BaseModel):
    step:int
    types:int
    amount:float	
    oldbalanceorig:float	
    newbalanceorig:float	
    oldbalancedest:float	
    newbalancedest:float	
    isflaggedfraud:float


@app.post('/predict')
def predict(data : fraudDetection):
                                                                                                                                                                                                                                
    features = np.array([[data.step, data.types, data.amount, data.oldbalanceorig, data.newbalanceorig, data.oldbalancedest, data.newbalancedest, data.isflaggedfraud]])
    model = joblib.load('credit_fraud.pkl')

    predictions = model.predict(features)
    if predictions == 1:
        return {"fraudulent"}
    elif predictions == 0:
        return {"not fraudulent"}

\end{lstlisting}

\section{app.py}
\begin{lstlisting}[language=Python]
import streamlit as st
import json
import requests as re

st.title("Credit Card Fraud Detection Web App")

st.image("image.png")

st.write("""
## About
Credit card fraud is a form of identity theft that involves an unauthorized taking of another's credit card information for the purpose of charging purchases to the account or removing funds from it.

**This Streamlit App utilizes a Machine Learning API in order to detect fraudulent credit card  based on the following criteria: hours, type of transaction, amount, balance before and after transaction etc.** 

The notebook, model and documentation(Dockerfiles, FastAPI script, Streamlit App script) are available on [GitHub.](https://github.com/Nneji123/Credit-Card-Fraud-Detection)        

**Developed by:**
- **Prakhar Bhardwaj**
- **Sourabh Kumar Verma**
""")""")


st.sidebar.header('Input Features of The Transaction')

sender_name = st.sidebar.text_input("""Input Sender ID""")
receiver_name = st.sidebar.text_input("""Input Receiver ID""")
step = st.sidebar.slider("""Number of Hours it took the Transaction to complete: """)
types = st.sidebar.subheader(f"""
                 Enter Type of Transfer Made:\n\n\n\n
                 0 for 'Cash In' Transaction\n 
                 1 for 'Cash Out' Transaction\n 
                 2 for 'Debit' Transaction\n
                 3 for 'Payment' Transaction\n  
                 4 for 'Transfer' Transaction\n""")
types = st.sidebar.selectbox("",(0,1,2,3,4))
x = ''
if types == 0:
    x = 'Cash in'
if types == 1:
    x = 'Cash Out'
if types == 2:
    x = 'Debit'
if types == 3:
    x = 'Payment'
if types == 4:
    x =  'Transfer'
    
amount = st.sidebar.number_input("Amount in $",min_value=0, max_value=110000)
oldbalanceorg = st.sidebar.number_input("""Original Balance Before Transaction was made""",min_value=0, max_value=110000)
newbalanceorg= st.sidebar.number_input("""New Balance After Transaction was made""",min_value=0, max_value=110000)
oldbalancedest= st.sidebar.number_input("""Old Balance""",min_value=0, max_value=110000)
newbalancedest= st.sidebar.number_input("""New Balance""",min_value=0, max_value=110000)
isflaggedfraud = st.sidebar.selectbox("""Specify if this was flagged as Fraud by your System: """,(0,1))


if st.button("Detection Result"):
    values = {
    "step": step,
    "types": types,
    "amount": amount,
    "oldbalanceorig": oldbalanceorg,
    "newbalanceorig": newbalanceorg,
    "oldbalancedest": oldbalancedest,
    "newbalancedest": newbalancedest,
    "isflaggedfraud": isflaggedfraud
    }


    st.write(f"""### These are the transaction details:\n
    Sender ID: {sender_name}
    Receiver ID: {receiver_name}
    1. Number of Hours it took to complete: {step}\n
    2. Type of Transaction: {x}\n
    3. Amount Sent: {amount}\n
    4. Sender Previous Balance Before Transaction: {oldbalanceorg}\n
    5. Sender New Balance After Transaction: {newbalanceorg}\n
    6. Recepient Balance Before Transaction: {oldbalancedest}\n
    7. Recepient Balance After Transaction: {newbalancedest}\n
    8. System Flag Fraud Status: {isflaggedfraud}
                """)

    res = re.post(f"http://backend.docker:8000/predict/",json=values)
    json_str = json.dumps(res.json())
    resp = json.loads(json_str)
    
    if sender_name=='' or receiver_name == '':
        st.write("Error! Please input Transaction ID or Names of Sender and Receiver!")
    else:
        st.write(f"""### The '{x}' transaction that took place between {sender_name} and {receiver_name} is {resp[0]}.""")










\end{lstlisting}

\section{train.py}
\begin{lstlisting}[language=Python]
import pandas as pd
import numpy as np
import os
import time
import joblib
import warnings
warnings.filterwarnings('ignore')

from sklearn.model_selection import train_test_split, StratifiedKFold
from sklearn.ensemble import (
    RandomForestClassifier,
    StackingClassifier,
    IsolationForest,
)
from sklearn.linear_model import LogisticRegression
from sklearn.preprocessing import StandardScaler
from sklearn.pipeline import Pipeline
from sklearn.metrics import (
    classification_report,
    f1_score,
    roc_auc_score,
    precision_recall_curve,
    confusion_matrix,
)
import xgboost as xgb
import lightgbm as lgb
from catboost import CatBoostClassifier

print("=" * 80)
print("ADVANCED FRAUD DETECTION MODEL TRAINING (v3.0 — Stacking + LGBM + CatBoost)")
print("=" * 80)

# ─────────────────────────────────────────────────────────────────────────────
# 1. LOAD DATA
# ─────────────────────────────────────────────────────────────────────────────
print("\n[1/6] Loading PaySim dataset...")
current_dir = os.path.dirname(os.path.abspath(__file__))
data_path = os.path.join(
    current_dir, "..", "data", "raw",
    "PS_20174392719_1491204439457_log_augmented.csv",
)

if not os.path.exists(data_path):
    print(f"Error: Dataset not found at {data_path}")
    exit(1)

data = pd.read_csv(data_path)
print(f"      Dataset loaded: {len(data):,} transactions")

# ─────────────────────────────────────────────────────────────────────────────
# 2. PREPROCESSING & FEATURE ENGINEERING
# ─────────────────────────────────────────────────────────────────────────────
print("\n[2/6] Preprocessing & Advanced Feature Engineering...")

data = data.fillna(0)

# ── Basic temporal features ──────────────────────────────────────────────────
data["hour"]        = data["step"] % 24
data["day_of_week"] = (data["step"] // 24) % 7
data["is_night"]    = (data["hour"] <= 6).astype(int)

# ── Error balance features (high signal for fraud) ───────────────────────────
data["errorBalanceOrig"] = data["newbalanceOrig"] + data["amount"] - data["oldbalanceOrg"]
data["errorBalanceDest"] = data["oldbalanceDest"] + data["amount"] - data["newbalanceDest"]

# ── Log-scaled amount ────────────────────────────────────────────────────────
data["amount_log"] = np.log1p(data["amount"])

# ── NEW: Ratio of amount to sender's original balance ────────────────────────
# How large is the transfer relative to what the sender had?
data["amount_to_balance_ratio"] = data["amount"] / (data["oldbalanceOrg"] + 1)

# ── NEW: Balance drain flag ──────────────────────────────────────────────────
# Did the sender's balance reach zero after this transaction?
data["balance_drain_flag"] = (data["newbalanceOrig"] == 0).astype(int)

# ── NEW: Destination balance increase ratio ───────────────────────────────────
# How much did the dest balance grow relative to the amount received?
data["dest_balance_increase_ratio"] = (
    (data["newbalanceDest"] - data["oldbalanceDest"]) / (data["amount"] + 1)
)

# ── NEW: Round-number amount flag ────────────────────────────────────────────
# Fraudsters often use psychologically round numbers
data["is_round_amount"] = (data["amount"] % 1000 == 0).astype(int)

# ── NEW: Destination-to-amount balance ratio ─────────────────────────────────
data["balance_dest_ratio"] = data["amount"] / (data["oldbalanceDest"] + 1)

print("      Basic features done.")

# ─────────────────────────────────────────────────────────────────────────────
# 3. SPLIT (before encoding to prevent leakage)
# ─────────────────────────────────────────────────────────────────────────────
print("\n[3/6] Splitting data (75% train / 25% test, stratified)...")

# Temporarily keep 'type' as string for encoding
cols_to_drop = ["nameOrig", "nameDest", "isFraud"]
feature_cols = [c for c in data.columns if c not in cols_to_drop]
X_raw = data[feature_cols].copy()
y     = data["isFraud"].copy()

# Ordinal encode 'type' temporarily so we can split — we re-encode after
type_mapping_str = {
    "PAYMENT": 3, "TRANSFER": 4, "CASH_OUT": 1, "DEBIT": 2, "CASH_IN": 0
}
if X_raw["type"].dtype == object:
    X_raw["type"] = X_raw["type"].map(type_mapping_str).fillna(3)

X_train_raw, X_test_raw, y_train, y_test = train_test_split(
    X_raw, y, test_size=0.25, random_state=27, stratify=y
)
print(f"      Train: {len(X_train_raw):,}  |  Test: {len(X_test_raw):,}")

# ─────────────────────────────────────────────────────────────────────────────
# TARGET ENCODING for 'type' — computed only on train, applied to both
# ─────────────────────────────────────────────────────────────────────────────
print("      Applying target encoding to 'type' feature (no leakage)...")

train_with_target = X_train_raw.copy()
train_with_target["isFraud"] = y_train.values

type_fraud_rate = (
    train_with_target.groupby("type")["isFraud"].mean().to_dict()
)
global_rate = y_train.mean()

X_train = X_train_raw.copy()
X_test  = X_test_raw.copy()
X_train["type_encoded"] = X_train["type"].map(type_fraud_rate).fillna(global_rate)
X_test["type_encoded"]  = X_test["type"].map(type_fraud_rate).fillna(global_rate)

print(f"      Type fraud rates: {type_fraud_rate}")

# ─────────────────────────────────────────────────────────────────────────────
# ISOLATION FOREST anomaly score feature
# ─────────────────────────────────────────────────────────────────────────────
print("      Fitting Isolation Forest for anomaly score feature...")
iso_features = [
    "amount", "amount_log", "amount_to_balance_ratio",
    "errorBalanceOrig", "errorBalanceDest",
    "balance_drain_flag", "dest_balance_increase_ratio",
]
iso_forest = IsolationForest(
    n_estimators=100, contamination=0.02, random_state=27, n_jobs=-1
)
iso_forest.fit(X_train[iso_features])

# decision_function: higher = more normal, lower = more anomalous → invert
X_train["isolation_score"] = -iso_forest.decision_function(X_train[iso_features])
X_test["isolation_score"]  = -iso_forest.decision_function(X_test[iso_features])
print("      Isolation Forest done.")

# Final feature list (drop raw 'type' — replaced by type_encoded)
FINAL_FEATURES = [c for c in X_train.columns if c != "type"]
X_train = X_train[FINAL_FEATURES]
X_test  = X_test[FINAL_FEATURES]

print(f"\n      Final feature set ({len(FINAL_FEATURES)} features):")
for f in FINAL_FEATURES:
    print(f"        • {f}")

# ─────────────────────────────────────────────────────────────────────────────
# 4. SMOTE (training set only)
# ─────────────────────────────────────────────────────────────────────────────
print("\n[4/6] Balancing training data with SMOTE...")
from imblearn.over_sampling import SMOTE
smote = SMOTE(random_state=27, k_neighbors=5)
X_train_bal, y_train_bal = smote.fit_resample(X_train, y_train)
print(f"      After SMOTE — Total: {len(X_train_bal):,}  |  Fraud: {(y_train_bal==1).sum():,}  |  Legit: {(y_train_bal==0).sum():,}")

# ─────────────────────────────────────────────────────────────────────────────
# 5. DEFINE BASE LEARNERS
# ─────────────────────────────────────────────────────────────────────────────
print("\n[5/6] Building Advanced Stacking Ensemble...")

print("   [A] Random Forest (200 trees)...")
rf = RandomForestClassifier(
    n_estimators=200,
    max_depth=20,
    min_samples_split=10,
    class_weight="balanced",
    n_jobs=-1,
    random_state=27,
)

print("   [B] XGBoost (200 estimators, tuned)...")
xgb_model = xgb.XGBClassifier(
    n_estimators=200,
    learning_rate=0.05,
    max_depth=6,
    subsample=0.8,
    colsample_bytree=0.8,
    scale_pos_weight=(y_train_bal == 0).sum() / (y_train_bal == 1).sum(),
    eval_metric="logloss",
    random_state=27,
    n_jobs=-1,
    verbosity=0,
)

print("   [C] LightGBM (leaf-wise tree growth, 300 estimators)...")
lgbm_model = lgb.LGBMClassifier(
    n_estimators=300,
    learning_rate=0.05,
    num_leaves=63,
    max_depth=-1,
    min_child_samples=20,
    subsample=0.8,
    colsample_bytree=0.8,
    class_weight="balanced",
    random_state=27,
    n_jobs=-1,
    verbose=-1,
)

print("   [D] CatBoost (ordered boosting, native categorical handling, 300 iters)...")
cat_model = CatBoostClassifier(
    iterations=300,
    learning_rate=0.05,
    depth=6,
    eval_metric="F1",
    random_seed=27,
    verbose=0,
    auto_class_weights="Balanced",
)

# ── Meta-learner (Level-1) ───────────────────────────────────────────────────
print("   [E] Meta-Learner: Logistic Regression (trained on OOF predictions)...")
meta_lr = Pipeline([
    ("scaler", StandardScaler()),
    ("lr", LogisticRegression(C=1.0, solver="lbfgs", max_iter=1000, random_state=27)),
])

# ── Stacking Classifier ──────────────────────────────────────────────────────
print("\n   Assembling StackingClassifier...")
stacking_clf = StackingClassifier(
    estimators=[
        ("rf",      rf),
        ("xgb",     xgb_model),
        ("lgbm",    lgbm_model),
        ("catboost", cat_model),
    ],
    final_estimator=meta_lr,
    cv=StratifiedKFold(n_splits=5, shuffle=True, random_state=27),
    stack_method="predict_proba",
    n_jobs=-1,
    passthrough=False,  # meta-learner sees only base model outputs
)

print("\n   Training Stacking Ensemble (5-fold CV for OOF meta-features)...")
print("   >> This will take 10-20 minutes -- please wait...\n")
t0 = time.time()
stacking_clf.fit(X_train_bal, y_train_bal)
train_time = time.time() - t0
print(f"\n   >> Training complete in {train_time/60:.1f} minutes.")

# ─────────────────────────────────────────────────────────────────────────────
# EVALUATE BASE LEARNERS
# ─────────────────────────────────────────────────────────────────────────────
print("\n[G] Evaluating base learners on held-out test set...")
base_learner_names = {"rf": "Random Forest", "xgb": "XGBoost", "lgbm": "LightGBM", "catboost": "CatBoost"}
for key, display_name in base_learner_names.items():
    clf = stacking_clf.named_estimators_[key]
    y_prob = clf.predict_proba(X_test)[:, 1]
    prec, rec, threshs = precision_recall_curve(y_test, y_prob)
    with np.errstate(divide="ignore", invalid="ignore"):
        fsc = np.where((prec + rec) == 0, 0, 2 * prec * rec / (prec + rec))
    best_idx = np.argmax(fsc[:-1])  # last element has no threshold
    best_thr = threshs[best_idx] if len(threshs) > best_idx else 0.5
    y_pred = (y_prob >= best_thr).astype(int)
    print(
        f"   [{display_name}] "
        f"F1: {f1_score(y_test, y_pred):.4f} | "
        f"AUC: {roc_auc_score(y_test, y_prob):.4f} | "
        f"Thresh: {best_thr:.4f}"
    )

# ─────────────────────────────────────────────────────────────────────────────
# EVALUATE STACKING ENSEMBLE + THRESHOLD TUNING
# ─────────────────────────────────────────────────────────────────────────────
print("\n[H] Evaluating full Stacking Ensemble...")
y_pred_proba = stacking_clf.predict_proba(X_test)[:, 1]

precision, recall, thresholds = precision_recall_curve(y_test, y_pred_proba)
with np.errstate(divide="ignore", invalid="ignore"):
    fscore = np.where(
        (precision + recall) == 0, 0,
        2 * precision * recall / (precision + recall)
    )
best_ix = np.argmax(fscore[:-1])
best_thresh = float(thresholds[best_ix]) if len(thresholds) > best_ix else 0.5

y_pred_final = (y_pred_proba >= best_thresh).astype(int)
final_f1  = f1_score(y_test, y_pred_final)
final_auc = roc_auc_score(y_test, y_pred_proba)

print(f"   Best Threshold : {best_thresh:.4f}")
print(f"   F1 Score       : {final_f1:.4f}")
print(f"   AUC-ROC        : {final_auc:.4f}")
print()
print(classification_report(y_test, y_pred_final, target_names=["Legitimate", "Fraud"]))

print("\n" + "=" * 80)
print(f"SELECTED: Advanced Stacking Ensemble (RF + XGB + LightGBM + CatBoost → LR)")
print(f"   F1 Score  : {final_f1:.4f}")
print(f"   AUC-ROC   : {final_auc:.4f}")
print(f"   Threshold : {best_thresh:.4f}")
print("=" * 80)

# ─────────────────────────────────────────────────────────────────────────────
# FEATURE IMPORTANCE (average of tree-based models)
# ─────────────────────────────────────────────────────────────────────────────
print("\n[5b] Computing ensemble feature importance...")
rf_imp   = stacking_clf.named_estimators_["rf"].feature_importances_
xgb_imp  = stacking_clf.named_estimators_["xgb"].feature_importances_
lgbm_imp = stacking_clf.named_estimators_["lgbm"].feature_importances_
# CatBoost importances
cat_imp  = stacking_clf.named_estimators_["catboost"].get_feature_importance()

# Normalise each to sum to 1
def norm(x): return x / (x.sum() + 1e-10)
avg_importance = (norm(rf_imp) + norm(xgb_imp) + norm(lgbm_imp) + norm(cat_imp)) / 4

print("\n   Top-10 Most Important Features:")
feat_imp_pairs = sorted(zip(FINAL_FEATURES, avg_importance), key=lambda x: x[1], reverse=True)
for feat, imp in feat_imp_pairs[:10]:
    bar = "█" * int(imp * 200)
    print(f"   {feat:<35} {imp:.4f}  {bar}")

# ─────────────────────────────────────────────────────────────────────────────
# 6. SAVE MODEL
# ─────────────────────────────────────────────────────────────────────────────
print("\n[6/6] Saving model artifacts...")
save_path = os.path.join(current_dir, "..", "models", "credit_fraud.pkl")

artifacts = {
    "model":               stacking_clf,
    "threshold":           best_thresh,
    "features":            FINAL_FEATURES,
    "feature_importances": avg_importance,
    "iso_forest":          iso_forest,
    "iso_features":        iso_features,
    "type_encoding":       type_fraud_rate,
    "global_fraud_rate":   global_rate,
    "model_version":       "v3.0-stacking",
    "base_learners":       ["RandomForest", "XGBoost", "LightGBM", "CatBoost"],
    "meta_learner":        "LogisticRegression",
}

joblib.dump(artifacts, save_path)
print(f"   ✅ Saved to: {save_path}")
print(f"   Model size: {os.path.getsize(save_path) / 1024 / 1024:.1f} MB")
print("\n   Restart the API server to load the new model. Streamlit will auto-detect it.\n")
print("=" * 80)
print("DONE — Advanced Stacking Ensemble ready for production.")
print("=" * 80)

\end{lstlisting}

\section{augment_data.py}
\begin{lstlisting}[language=Python]
import pandas as pd
import numpy as np
import os
import random

print("="*80)
print("REALISTIC FRAUD DATA AUGMENTATION")
print("="*80)

# 1. Load Data
current_dir = os.path.dirname(os.path.abspath(__file__))
data_path = os.path.join(current_dir, '..', 'data', 'raw', 'PS_20174392719_1491204439457_log.csv')
augmented_data_path = os.path.join(current_dir, '..', 'data', 'raw', 'PS_20174392719_1491204439457_log_augmented.csv')

print(f"Loading data from {data_path}...")
df = pd.read_csv(data_path)
print(f"Original dataset size: {len(df)} rows")

# 2. Define targets for augmentation
target_types = ['PAYMENT', 'CASH_IN', 'DEBIT']
samples_per_type = 250
new_fraud_rows = []

print("\nSynthesizing realistic fraud cases (40/40/20 mix)...")
np.random.seed(42)
random.seed(42)

for t in target_types:
    # Get non-fraud transactions of this type
    eligible_rows = df[(df['type'] == t) & (df['isFraud'] == 0)]
    
    # Randomly sample cases
    sampled = eligible_rows.sample(n=samples_per_type, random_state=42).copy()
    
    # Set base fraud class
    sampled['isFraud'] = 1
    
    # Split into 3 buckets: 40% (100), 40% (100), 20% (50)
    bucket1_idx = sampled.index[:100]
    bucket2_idx = sampled.index[100:200]
    bucket3_idx = sampled.index[200:250]
    
    # 🟢 Type 1: Partial draining (40% - indices 0 to 99)
    # Large suspicious withdrawal, not obvious fraud.
    subset1 = sampled.loc[bucket1_idx].copy()
    random_factors = np.random.uniform(0.6, 0.95, size=len(subset1))
    subset1['amount'] = subset1['oldbalanceOrg'] * random_factors
    subset1['newbalanceOrig'] = subset1['oldbalanceOrg'] - subset1['amount']
    
    # 🟡 Type 2: Balance inconsistency (40% - indices 100 to 199)
    # amount is deducted, but balance doesn't reflect it
    subset2 = sampled.loc[bucket2_idx].copy()
    subset2['amount'] = np.random.uniform(1000, 15000, size=len(subset2))
    subset2['newbalanceOrig'] = subset2['oldbalanceOrg']  # No deduction applied
    
    # 🔴 Type 3: Impossible balance (20% - indices 200 to 249)
    # Balance increases after spending!
    subset3 = sampled.loc[bucket3_idx].copy()
    subset3['amount'] = subset3['oldbalanceOrg'] * 0.8
    subset3['newbalanceOrig'] = subset3['oldbalanceOrg'] + 2000
    
    # Combine back
    synthetic_batch = pd.concat([subset1, subset2, subset3])
    new_fraud_rows.append(synthetic_batch)
    print(f"  - Generated {len(synthetic_batch)} synthetic fraud rows for {t}")

# 3. Combine and shuffle
synthetic_df = pd.concat(new_fraud_rows, ignore_index=True)
augmented_df = pd.concat([df, synthetic_df], ignore_index=True)

# Shuffle the dataframe to distribute the continuous block of new records evenly
augmented_df = augmented_df.sample(frac=1, random_state=42).reset_index(drop=True)

print(f"\nAugmented dataset size: {len(augmented_df)} rows")

# Verify the fraud distribution
print("\nNew Fraud Distribution by Type:")
print(augmented_df.groupby('type')['isFraud'].sum())

# 4. Save the new dataset
print(f"\nSaving augmented dataset to {augmented_data_path}...")
augmented_df.to_csv(augmented_data_path, index=False)
print("Done!")

\end{lstlisting}

\section{download_animations.py}
\begin{lstlisting}[language=Python]
import requests
import json

def download_lottie(url, filename):
    response = requests.get(url)
    if response.status_code == 200:
        with open(filename, 'w') as f:
            json.dump(response.json(), f)
        print(f"Downloaded {filename}")
    else:
        print(f"Failed to download {filename}")

# Security Shield Animation (Using Scan Animation as per user request)
download_lottie("https://assets8.lottiefiles.com/packages/lf20_qp1q7mct.json", "shield.json")
# Scanning Animation
download_lottie("https://assets8.lottiefiles.com/packages/lf20_qp1q7mct.json", "scan.json")

\end{lstlisting}

\section{generate_doc.py}
\begin{lstlisting}[language=Python]
import os
import docx
from docx.shared import Inches, Pt, RGBColor
from docx.enum.text import WD_ALIGN_PARAGRAPH
from docx.enum.style import WD_STYLE_TYPE

def create_ieee_word_doc():
    doc = docx.Document()
    
    # Configure styles
    style = doc.styles['Normal']
    font = style.font
    font.name = 'Times New Roman'
    font.size = Pt(10)

    # Title
    title = doc.add_heading('FraudGuard AI: A Heterogeneous Ensemble Approach to Real-Time Financial Fraud Detection', level=0)
    title.alignment = WD_ALIGN_PARAGRAPH.CENTER
    title_font = title.style.font
    title_font.name = 'Times New Roman'
    title_font.size = Pt(24)
    # Removing bold from title standard style as it's already a heading but can force if needed
    
    doc.add_paragraph('\n')

    # Abstract & Index Terms
    p = doc.add_paragraph()
    runner = p.add_run('Abstract— ')
    runner.bold = True
    p.add_run('The rapid digitalization of financial services has precipitated a corresponding escalation in sophisticated fraudulent activities. Traditional rule-based systems are increasingly incapable of detecting anomalous transaction patterns in real-time without hindering legitimate user experiences. This paper presents FraudGuard AI, an advanced, robust machine learning pipeline designed to detect financial fraud with high precision. Our proposed methodology leverages a heterogeneous soft-voting ensemble mechanism comprising Random Forest (RF), eXtreme Gradient Boosting (XGBoost), AdaBoost, Logistic Regression (LR), and a Deep Neural Network (MLP). By engineering domain-specific features related to ledger discrepancies, the system effectively isolates malicious account takeovers. Evaluated on the widely recognized PaySim dataset (simulating mobile money transactions), the ensemble model achieves an exceptional F1-Score of 0.989, demonstrating superior capability in balancing precision and recall within a highly imbalanced dataset context. Furthermore, we detail the deployment architecture utilizing FastAPI and Streamlit to facilitate real-time inference and model explainability via SHAP values.')
    p.alignment = WD_ALIGN_PARAGRAPH.JUSTIFY

    p2 = doc.add_paragraph()
    runner2 = p2.add_run('Index Terms— ')
    runner2.bold = True
    p2.add_run('Machine Learning, Fraud Detection, Ensemble Methods, Feature Engineering, Real-Time Systems, SHAP Explainability.\n')
    
    # Section I. Introduction
    sec1 = doc.add_heading('I. INTRODUCTION', level=1)
    sec1.alignment = WD_ALIGN_PARAGRAPH.CENTER
    doc.add_paragraph('The proliferation of electronic payment systems has fundamentally transformed the global economic landscape. However, this convenience is accompanied by the persistent threat of financial fraud. In particular, mobile money transfers and credit card transactions are attractive targets for cybercriminals. Fraud detection systems must operate under stringent constraints: they must analyze immense volumes of transaction data in milliseconds, minimize false positives to avoid customer friction, and continuously adapt to novel adversarial strategies.').alignment = WD_ALIGN_PARAGRAPH.JUSTIFY
    doc.add_paragraph('Traditional fraud detection relies heavily on rigid, rule-based systems defined by domain experts. These systems suffer from high maintenance overhead and struggle to generalize to unseen fraud patterns. Consequently, the industry has shifted towards Machine Learning (ML) techniques. While standalone models like Random Forests and Support Vector Machines yield promising results, the complex, non-linear, and highly imbalanced nature of transactional data (where fraudulent transactions represent less than 0.2% of the dataset) necessitates more advanced architectures.').alignment = WD_ALIGN_PARAGRAPH.JUSTIFY
    doc.add_paragraph('In this research, we introduce FraudGuard AI, a comprehensive ecosystem for predictive fraud monitoring. The primary contributions of this paper are:').alignment = WD_ALIGN_PARAGRAPH.JUSTIFY
    # List
    doc.add_paragraph('1) The formulation of domain-engineered features (e.g., errorBalanceOrig) that expose underlying mathematical discrepancies.', style='List Number')
    doc.add_paragraph('2) The development of a Penta-model heterogeneous ensemble (Soft Voting Classifier) comprising Bagging, Boosting, Linear, and Deep Learning algorithms.', style='List Number')
    doc.add_paragraph('3) The implementation of a real-time explainability framework utilizing SHAP to provide transparent insights into the model\'s decision-making process.', style='List Number')

    # Section II. Related Work
    sec2 = doc.add_heading('II. RELATED WORK', level=1)
    sec2.alignment = WD_ALIGN_PARAGRAPH.CENTER
    doc.add_paragraph('Recent literature extensively evaluates Machine Learning for fraud detection. Dal Pozzolo et al. highlighted the challenges of concept drift and class imbalance in credit card fraud. Ensemble methods are frequently proposed to mitigate these issues; for instance, boosting algorithms like XGBoost have shown state-of-the-art performance in numerous tabular data competitions.').alignment = WD_ALIGN_PARAGRAPH.JUSTIFY
    doc.add_paragraph('However, many existing studies focus solely on offline metrics without addressing the engineering challenges of real-time deployment and interpretability. Our approach extends existing paradigms by integrating model threshold optimization with a production-ready application layer that explicitly visualizes the contribution of specific features for end-users.').alignment = WD_ALIGN_PARAGRAPH.JUSTIFY

    # Section III. Methodology
    sec3 = doc.add_heading('III. METHODOLOGY', level=1)
    sec3.alignment = WD_ALIGN_PARAGRAPH.CENTER
    doc.add_paragraph('Our methodology encompasses data curation, extensive feature engineering, and the training of a soft-voting ensemble classifier.').alignment = WD_ALIGN_PARAGRAPH.JUSTIFY
    
    doc.add_heading('A. Dataset Description', level=2)
    doc.add_paragraph('The model is trained and evaluated using the PaySim dataset, a synthetically generated but highly realistic dataset based on one month of financial logs from a mobile money service implemented in an African country. For our experimental setup, a subset of 1,048,575 transactions was utilized. The class imbalance ratio is ~0.11% fraud prevalence.').alignment = WD_ALIGN_PARAGRAPH.JUSTIFY

    doc.add_heading('B. Preprocessing and Feature Engineering', level=2)
    doc.add_paragraph('Given the extreme class imbalance, raw transactional data is insufficient. We applied mapping to categorical variables and engineered several critical features indicative of anomalous behavior. Crucially, we implemented ledger discrepancy features such as errorBalanceOrig, calculated as newbalanceOrig + amount - oldbalanceOrg. In legitimate transfers, this should equal zero. Deviations are absolute mathematical indicators of forced balance modifications or system anomalies.').alignment = WD_ALIGN_PARAGRAPH.JUSTIFY

    doc.add_heading('C. The Heterogeneous Ensemble Architecture', level=2)
    doc.add_paragraph('To maximize predictive robustness, we designed a soft-voting ensemble combining five algorithms: Random Forest (Bagging), XGBoost (Boosting), AdaBoost, Logistic Regression (Linear), and a Deep Neural Network (MLP). The Soft-Voting Classifier averages the predicted probabilities of the models. The final classification boundary (Threshold) is dynamically tuned to maximize the F1-Score via Precision-Recall curve analysis.').alignment = WD_ALIGN_PARAGRAPH.JUSTIFY

    # Section IV. Experimental Setup
    sec4 = doc.add_heading('IV. EXPERIMENTAL SETUP AND RESULTS', level=1)
    sec4.alignment = WD_ALIGN_PARAGRAPH.CENTER
    doc.add_paragraph('The dataset was stratified and split into a 75% training set (786,431 records) and a 25% testing set (262,144 records).').alignment = WD_ALIGN_PARAGRAPH.JUSTIFY
    
    doc.add_heading('A. Model Evaluation Metrics', level=2)
    doc.add_paragraph('Given the 0.11% minority class, Accuracy is a deceptive metric. We utilized the F1-Score (harmonic mean of Precision and Recall) and the Area Under the Receiver Operating Characteristic Curve (ROC-AUC).')
    
    doc.add_heading('B. Results and Comparison', level=2)
    doc.add_paragraph('The empirical results demonstrate the efficacy of the ensemble approach. The table and figure below illustrate the performance of individual models against the unified Penta-Ensemble.').alignment = WD_ALIGN_PARAGRAPH.JUSTIFY

    # Add Table
    table = doc.add_table(rows=1, cols=3)
    table.style = 'Light Shading Accent 1'
    hdr_cells = table.rows[0].cells
    hdr_cells[0].text = 'Model Architecture'
    hdr_cells[1].text = 'F1-Score'
    hdr_cells[2].text = 'ROC-AUC'
    
    data = [
        ('Logistic Regression', '0.6192', '0.9563'),
        ('AdaBoost', '0.9912', '0.9926'),
        ('Deep Neural Network', '0.9820', '0.9880'),
        ('XGBoost', '0.9900', '0.9950'),
        ('Random Forest', '0.9850', '0.9910'),
        ('Penta-Ensemble', '0.9894', '0.9960')
    ]
    for m, f, a in data:
        row_cells = table.add_row().cells
        row_cells[0].text = m
        row_cells[1].text = f
        row_cells[2].text = a

    doc.add_paragraph('\n')
    
    # Add Image
    try:
        if os.path.exists('docs/assets/model_comparison.png'):
            doc.add_picture('docs/assets/model_comparison.png', width=Inches(5.0))
            last_paragraph = doc.paragraphs[-1]
            last_paragraph.alignment = WD_ALIGN_PARAGRAPH.CENTER
            doc.add_paragraph('Fig 1. Performance Comparison of Evaluated Models').alignment = WD_ALIGN_PARAGRAPH.CENTER
    except Exception as e:
        print(f"Could not insert image: {e}")

    # Section V. Deployment
    sec5 = doc.add_heading('V. SYSTEM DEPLOYMENT AND EXPLAINABILITY', level=1)
    sec5.alignment = WD_ALIGN_PARAGRAPH.CENTER
    doc.add_paragraph('A critical barrier to deploying ML in finance is the "black-box" nature of complex models. FraudGuard AI addresses this via a decoupled microservices architecture utilizing FastAPI for backend inference and Streamlit for the frontend dashboard. For every flagged transaction, the system computes and visualizes SHAP impact values immediately, rendering a waterfall chart that explicitly dictates why a transaction was deemed fraudulent (e.g., displaying errorBalanceOrig (+5.26) to denote massive deviation).').alignment = WD_ALIGN_PARAGRAPH.JUSTIFY

    # Section VI. Conclusion
    sec6 = doc.add_heading('VI. CONCLUSION', level=1)
    sec6.alignment = WD_ALIGN_PARAGRAPH.CENTER
    doc.add_paragraph('This paper introduced an end-to-end framework for detecting financial transaction fraud using a heterogeneous soft-voting ensemble comprising traditional and deep learning architectures. By effectively augmenting the dataset with mathematical ledger validation features, the system achieves an F1-Score of 0.989. The subsequent deployment of the model ensures that FraudGuard AI is not merely a theoretical construct, but a functional, auditable, and resilient production-ready application.').alignment = WD_ALIGN_PARAGRAPH.JUSTIFY

    # Save
    os.makedirs('docs', exist_ok=True)
    doc.save('docs/FraudGuard_IEEE_Paper.docx')
    print("Word document generated at docs/FraudGuard_IEEE_Paper.docx")

if __name__ == '__main__':
    create_ieee_word_doc()

\end{lstlisting}

\section{generate_plots.py}
\begin{lstlisting}[language=Python]
import pandas as pd
import numpy as np
from sklearn.model_selection import train_test_split
from sklearn.metrics import f1_score, roc_auc_score, precision_recall_curve, confusion_matrix
import joblib
import os
import matplotlib.pyplot as plt
import seaborn as sns

def generate_plots():
    print("Generating plots...")
    current_dir = os.path.dirname(os.path.abspath(__file__))
    base_dir = os.path.abspath(os.path.join(current_dir, '..'))
    data_path = os.path.join(base_dir, 'data', 'raw', 'PS_20174392719_1491204439457_log_augmented.csv')
    model_path = os.path.join(base_dir, 'models', 'credit_fraud.pkl')
    output_dir = os.path.join(base_dir, 'docs', 'Paper Review PDF')
    
    os.makedirs(output_dir, exist_ok=True)
    
    print("Loading data...")
    data = pd.read_csv(data_path)
    
    type_mapping = {
        "PAYMENT": 3,
        "TRANSFER": 4,
        "CASH_OUT": 1,
        "DEBIT": 2,
        "CASH_IN": 0
    }
    data["type"] = data["type"].map(type_mapping)
    data = data.fillna(0)
    
    data["hour"] = data["step"] % 24
    data["errorBalanceOrig"] = data["newbalanceOrig"] + data["amount"] - data["oldbalanceOrg"]
    data["errorBalanceDest"] = data["oldbalanceDest"] + data["amount"] - data["newbalanceDest"]
    
    data["day_of_week"] = (data["step"] // 24) % 7
    data["is_night"] = data["hour"].apply(lambda x: 1 if x >= 0 and x <= 6 else 0)
    data["amount_log"] = np.log1p(data["amount"])
    data["balance_dest_ratio"] = data["amount"] / (data["oldbalanceDest"] + 1)
    
    cols_to_drop = ['nameOrig', 'nameDest', 'isFraud']
    X = data.drop([c for c in cols_to_drop if c in data.columns], axis=1)
    y = data["isFraud"]
    
    print("Splitting data...")
    X_train, X_test, y_train, y_test = train_test_split(
        X, y, test_size=0.25, random_state=27, stratify=y
    )
    
    print("Loading model...")
    artifacts = joblib.load(model_path)
    voting_clf = artifacts['model']
    best_thresh_voting = artifacts['threshold']
    
    print("Predicting ensemble...")
    y_pred_proba_voting = voting_clf.predict_proba(X_test)[:, 1]
    y_pred_voting = (y_pred_proba_voting >= best_thresh_voting).astype(int)
    
    print("Generating confusion matrix...")
    cm = confusion_matrix(y_test, y_pred_voting)
    plt.figure(figsize=(8, 6))
    sns.heatmap(cm, annot=True, fmt="d", cmap="Blues", xticklabels=["Legitimate", "Fraudulent"], yticklabels=["Legitimate", "Fraudulent"])
    plt.title("Confusion Matrix - Penta-Ensemble Model")
    plt.xlabel("Predicted")
    plt.ylabel("Actual")
    cm_path = os.path.join(output_dir, 'confusion_matrix.png')
    plt.savefig(cm_path, dpi=300, bbox_inches='tight')
    plt.close()
    print(f"Saved {cm_path}")
    
    print("Evaluating individual models...")
    print(f"Available estimators in saved model: {list(voting_clf.named_estimators_.keys())}")
    models_to_eval = []
    for name, clf in voting_clf.named_estimators_.items():
        if clf is not None and clf != 'drop':
            models_to_eval.append((name.upper(), clf))

    
    results = []
    for name, clf in models_to_eval:
        print(f"Predicting {name}...")
        y_pred_proba = clf.predict_proba(X_test)[:, 1]
        prec, rec, threshs = precision_recall_curve(y_test, y_pred_proba)
        with np.errstate(divide='ignore', invalid='ignore'):
            fsc = np.where((prec + rec) == 0, 0, (2 * prec * rec) / (prec + rec))
        best_idx = np.argmax(fsc)
        best_thr = threshs[best_idx]
        y_pred = (y_pred_proba >= best_thr).astype(int)
        f1 = f1_score(y_test, y_pred)
        auc = roc_auc_score(y_test, y_pred_proba)
        results.append({'Model': name, 'F1-Score': f1, 'AUC': auc})
    
    voting_f1 = f1_score(y_test, y_pred_voting)
    voting_auc = roc_auc_score(y_test, y_pred_proba_voting)
    results.append({'Model': 'Penta-Ensemble', 'F1-Score': voting_f1, 'AUC': voting_auc})
    
    df_res = pd.DataFrame(results)
    
    print("Generating model comparison plot...")
    plt.figure(figsize=(12, 7))
    df_melted = df_res.melt(id_vars='Model', var_name='Metric', value_name='Score')
    sns.barplot(x='Model', y='Score', hue='Metric', data=df_melted, palette='viridis')
    
    # Adjust y-limits dynamically
    min_score = df_melted['Score'].min()
    plt.ylim(max(0.0, min_score - 0.1), 1.0)
    
    plt.title("Model Comparison: F1-Score & AUC", fontsize=14)
    plt.xticks(rotation=45, ha='right', fontsize=12)
    plt.legend(loc='lower right', fontsize=12)
    plt.tight_layout()
    comp_path = os.path.join(output_dir, 'model_comparison.png')
    plt.savefig(comp_path, dpi=300)
    plt.close()
    print(f"Saved {comp_path}")
    
if __name__ == '__main__':
    generate_plots()

\end{lstlisting}

\section{app.py}
\begin{lstlisting}[language=Python]
import streamlit as st
import json
import requests as re
from streamlit_lottie import st_lottie
import time
import joblib
import pandas as pd
import numpy as np
import matplotlib.pyplot as plt
import altair as alt
import os

# Page Config
st.set_page_config(
    page_title="FraudGuard AI",
    page_icon="🛡️",
    layout="wide",
    initial_sidebar_state="expanded"
)

# Session State Initialization
from datetime import datetime, timedelta

# --- Demo prefill defaults ---
if 'demo_step' not in st.session_state:
    st.session_state.demo_step = 1
if 'demo_amount' not in st.session_state:
    st.session_state.demo_amount = 0.0
if 'demo_type_idx' not in st.session_state:
    st.session_state.demo_type_idx = 3  # Payment
if 'demo_oldbalanceorg' not in st.session_state:
    st.session_state.demo_oldbalanceorg = 0.0
if 'demo_newbalanceorg' not in st.session_state:
    st.session_state.demo_newbalanceorg = 0.0
if 'demo_oldbalancedest' not in st.session_state:
    st.session_state.demo_oldbalancedest = 0.0
if 'demo_newbalancedest' not in st.session_state:
    st.session_state.demo_newbalancedest = 0.0
if 'demo_sender' not in st.session_state:
    st.session_state.demo_sender = ""
if 'demo_receiver' not in st.session_state:
    st.session_state.demo_receiver = ""

if 'history' not in st.session_state:
    now = datetime.now()
    st.session_state.history = [
        {"Timestamp": (now - timedelta(minutes=20)).strftime("%H:%M:%S"), "Source": "C1239821", "Destination": "M928371", "Amount": 450.00, "Risk Score": 0.02, "Status": "Safe"},
        {"Timestamp": (now - timedelta(minutes=16)).strftime("%H:%M:%S"), "Source": "C9827312", "Destination": "M128371", "Amount": 12500.00, "Risk Score": 0.15, "Status": "Safe"},
        {"Timestamp": (now - timedelta(minutes=14)).strftime("%H:%M:%S"), "Source": "C6271823", "Destination": "C721638", "Amount": 250000.00, "Risk Score": 0.89, "Status": "Fraud"},
        {"Timestamp": (now - timedelta(minutes=10)).strftime("%H:%M:%S"), "Source": "C1283712", "Destination": "M827162", "Amount": 50.00, "Risk Score": 0.005, "Status": "Safe"},
        {"Timestamp": (now - timedelta(minutes=7)).strftime("%H:%M:%S"), "Source": "C7261523", "Destination": "C261524", "Amount": 9800.00, "Risk Score": 0.04, "Status": "Safe"},
        {"Timestamp": (now - timedelta(minutes=3)).strftime("%H:%M:%S"), "Source": "C1827361", "Destination": "M928172", "Amount": 180000.00, "Risk Score": 0.76, "Status": "Fraud"},
        {"Timestamp": (now - timedelta(seconds=30)).strftime("%H:%M:%S"), "Source": "C9182736", "Destination": "M129381", "Amount": 120.50, "Risk Score": 0.01, "Status": "Safe"},
    ]

# Load Lottie Animation
@st.cache_resource
def load_lottiefile(filename: str):
    try:
        # Construct absolute path to assets
        current_dir = os.path.dirname(os.path.abspath(__file__))
        filepath = os.path.join(current_dir, 'assets', filename)
        with open(filepath, "r") as f:
            return json.load(f)
    except FileNotFoundError:
        return None

# Custom CSS for "Pro" Look
st.markdown("""
<style>
    /* Google Fonts */
    @import url('https://fonts.googleapis.com/css2?family=Roboto:wght@300;400;700&display=swap');
    
    html, body, [class*="css"] {
        font-family: 'Roboto', sans-serif;
    }
    
    /* Header Styling */
    .main-header {
        font-size: 2.5rem;
        font-weight: 700;
        background: -webkit-linear-gradient(#fca311, #e5e5e5);
        -webkit-background-clip: text;
        -webkit-text-fill-color: transparent;
        margin-bottom: 0px;
    }
    
    /* Card Styling */
    .stCard {
        background-color: #1e1e1e;
        padding: 20px;
        border-radius: 10px;
        box-shadow: 0 4px 6px rgba(0,0,0,0.3);
        border: 1px solid #333;
    }
    
    /* Metric Styling */
    div[data-testid="stMetricValue"] {
        font-size: 1.8rem;
    }
    
    /* Button Styling */
    .stButton>button {
        width: 100%;
        background-color: #fca311;
        color: black;
        font-weight: bold;
        border: none;
        padding: 10px;
        border-radius: 5px;
        transition: all 0.3s ease;
    }
    .stButton>button:hover {
        background-color: #e59400;
        color: white;
    }
</style>
""", unsafe_allow_html=True)

# Sidebar
lottie_shield = load_lottiefile("shield.json")
with st.sidebar:
    if lottie_shield:
        st_lottie(lottie_shield, height=150, key="shield")
    else:
        current_dir = os.path.dirname(os.path.abspath(__file__))
        img_path = os.path.join(current_dir, 'assets', 'image.png')
        if os.path.exists(img_path):
            st.image(img_path, use_column_width="auto")

    st.title("🛡️ FraudGuard AI")

    # --- API Health & Stats Widget ---
    try:
        api_url = os.getenv("API_URL", "http://127.0.0.1:8000")
        h_res = re.get(f"{api_url}/health", timeout=1)
        s_res = re.get(f"{api_url}/stats", timeout=1)
        if h_res.status_code == 200 and s_res.status_code == 200:
            h = h_res.json()
            s = s_res.json()
            st.markdown("### System Status: 🟢 Online")
            st.markdown(
                f"""
                <div style='background:#1a1a2e;border-radius:10px;padding:12px 14px;border:1px solid #333;margin-bottom:8px'>
                    <div style='color:#aaa;font-size:0.75rem;text-transform:uppercase;letter-spacing:1px;margin-bottom:8px'>⚡ API Live Metrics</div>
                    <div style='display:flex;justify-content:space-between;margin-bottom:6px'>
                        <span style='color:#ccc;font-size:0.85rem'>Uptime</span>
                        <span style='color:#fca311;font-weight:600;font-size:0.85rem'>{h['uptime']}</span>
                    </div>
                    <div style='display:flex;justify-content:space-between;margin-bottom:6px'>
                        <span style='color:#ccc;font-size:0.85rem'>Total Scanned</span>
                        <span style='color:#fff;font-weight:600;font-size:0.85rem'>{s['total_predictions']}</span>
                    </div>
                    <div style='display:flex;justify-content:space-between;margin-bottom:6px'>
                        <span style='color:#ccc;font-size:0.85rem'>Fraud Detected</span>
                        <span style='color:#ff4b4b;font-weight:600;font-size:0.85rem'>{s['fraud_count']}</span>
                    </div>
                    <div style='display:flex;justify-content:space-between'>
                        <span style='color:#ccc;font-size:0.85rem'>Fraud Rate</span>
                        <span style='color:#{'ff4b4b' if s['fraud_rate_pct'] > 20 else 'fca311'};font-weight:600;font-size:0.85rem'>{s['fraud_rate_pct']}%</span>
                    </div>
                </div>
                """,
                unsafe_allow_html=True
            )
        else:
            st.markdown("### System Status: 🟡 Degraded")
    except Exception:
        st.markdown("### System Status: 🔴 API Offline")

    # --- Live Threat Feed ---
    st.markdown("---")
    st.markdown("### 🔴 Live Threat Feed")
    recent = list(reversed(st.session_state.history[-5:])) if st.session_state.history else []
    if recent:
        for tx in recent:
            is_fraud = tx['Status'] == 'Fraud'
            border_color = '#ff4b4b' if is_fraud else '#22c55e'
            icon = '🚨' if is_fraud else '✅'
            risk_pct = f"{tx['Risk Score']:.0%}" if isinstance(tx['Risk Score'], float) else tx['Risk Score']
            amount_fmt = f"${tx['Amount']:,.0f}" if isinstance(tx['Amount'], (int, float)) else tx['Amount']
            st.markdown(
                f"""
                <div style='background:#111827;border-left:3px solid {border_color};border-radius:6px;
                            padding:8px 10px;margin-bottom:6px;'>
                    <div style='display:flex;justify-content:space-between;align-items:center'>
                        <span style='font-size:0.8rem;color:#9ca3af'>{tx['Timestamp']}</span>
                        <span style='font-size:0.75rem;color:{border_color};font-weight:700'>{icon} {tx['Status'].upper()}</span>
                    </div>
                    <div style='font-size:0.78rem;color:#d1d5db;margin-top:3px'>
                        <b style='color:#f9fafb'>{str(tx['Source'])[:10]}</b> → <b style='color:#f9fafb'>{str(tx['Destination'])[:10]}</b>
                    </div>
                    <div style='display:flex;justify-content:space-between;margin-top:3px'>
                        <span style='font-size:0.75rem;color:#fca311'>{amount_fmt}</span>
                        <span style='font-size:0.75rem;color:#9ca3af'>Risk: {risk_pct}</span>
                    </div>
                </div>
                """,
                unsafe_allow_html=True
            )
    else:
        st.caption("No transactions yet this session.")

    st.markdown("---")
    st.info("""
    **Real-Time Fraud Detection System**
    Detects anomalies in financial transactions using a Heterogeneous Penta-Ensemble Framework.

    **Developed by:**
    - **Prakhar Bhardwaj**
    - **Sourabh Kumar Verma**
    - **Satyam Kumar**
    - **Vivek Singh Bisht**
    """)

# Main Content
# Dashboard KPIs
st.markdown("### 📊 Live Command Center")
kpi1, kpi2, kpi3, kpi4 = st.columns(4)

# Calculate Metrics
total_scans = len(st.session_state.history)
fraud_detected = len([x for x in st.session_state.history if x['Status'] == 'Fraud'])
fraud_value = sum([x['Amount'] for x in st.session_state.history if x['Status'] == 'Fraud'])
avg_risk = sum([x['Risk Score'] for x in st.session_state.history]) / total_scans if total_scans > 0 else 0

kpi1.metric("Total Scans", total_scans, delta=f"+{total_scans} (Session)")
kpi2.metric("Fraud Detected", fraud_detected, delta="CRITICAL" if fraud_detected > 0 else "Normal", delta_color="inverse")
kpi3.metric("Blocked Volume", f"${fraud_value:,.2f}", delta="Saved Value")
kpi4.metric("Avg Risk Level", f"{avg_risk:.1%}", delta="Network Avg")

st.divider()

c1, c2 = st.columns([3, 1])
with c1:
    st.markdown('<p class="main-header">🛡️ Secure Transaction Analysis</p>', unsafe_allow_html=True)
    st.markdown("Enter transaction details below to scan for potential security threats.")

st.divider()

tab1, tab2, tab3 = st.tabs(["📝 Single Transaction", "📂 Batch Upload", "🧠 Model Insights"])

with tab3:
    st.subheader("🎓 Project Architecture & Model Explainability")

    # Architecture Diagram
    st.markdown("#### 1. System Architecture")
    st.markdown("""
    This project uses a microservices-inspired architecture:
    - **Frontend**: Streamlit (User Interface)
    - **Backend**: FastAPI (Prediction Logic)
    - **Model**: Heterogeneous Penta-Ensemble Framework (RF, XGB, Ada, LR, MLP)
    """)

    st.graphviz_chart("""
    digraph G {
        rankdir=LR;
        node [shape=box, style=filled, fillcolor=lightblue];
        User [shape=ellipse, fillcolor=lightgrey];
        Streamlit [label="Frontend UI\\n(Streamlit)"];
        FastAPI [label="Backend API\\n(FastAPI)"];
        Model [label="ML Model\\n(Penta-Ensemble)", fillcolor=lightgreen];

        User -> Streamlit [label="Interacts"];
        Streamlit -> FastAPI [label="Sends Data (JSON)"];
        FastAPI -> Model [label="Features"];
        Model -> FastAPI [label="Prediction"];
        FastAPI -> Streamlit [label="Response"];
    }
    """)

    st.divider()

    # --- Transaction Network Graph ---
    st.markdown("#### 2. Transaction Network Graph")
    st.info("Each bubble is a transaction from your session. **Size** = Amount, **Colour** = Risk Score (red = high risk). Hover for details.")

    if len(st.session_state.history) >= 2:
        net_df = pd.DataFrame(st.session_state.history).copy()
        net_df["Order"] = range(len(net_df))
        # Jitter X slightly so overlapping nodes are visible
        np.random.seed(42)
        net_df["jitter_x"] = net_df["Order"] + np.random.uniform(-0.3, 0.3, len(net_df))
        net_df["jitter_y"] = np.random.uniform(0.2, 0.8, len(net_df))
        net_df["AmountSafe"] = net_df["Amount"].apply(lambda x: x if isinstance(x, (int, float)) else 0)
        net_df["RiskSafe"] = net_df["Risk Score"].apply(lambda x: x if isinstance(x, (int, float)) else 0)
        # Normalise bubble size
        max_amt = net_df["AmountSafe"].max() or 1
        net_df["BubbleSize"] = ((net_df["AmountSafe"] / max_amt) * 900 + 100).clip(upper=1200)

        network_chart = alt.Chart(net_df).mark_circle(opacity=0.85, stroke="#1e1e2e", strokeWidth=1).encode(
            x=alt.X("jitter_x:Q", axis=None, scale=alt.Scale(zero=False)),
            y=alt.Y("jitter_y:Q", axis=None, scale=alt.Scale(domain=[0, 1])),
            size=alt.Size("BubbleSize:Q", legend=None, scale=alt.Scale(range=[80, 1200])),
            color=alt.Color(
                "RiskSafe:Q",
                scale=alt.Scale(domain=[0, 0.5, 1.0], range=["#22c55e", "#fca311", "#ff4b4b"]),
                legend=alt.Legend(title="Risk Score")
            ),
            tooltip=[
                alt.Tooltip("Timestamp", title="Time"),
                alt.Tooltip("Source", title="Sender"),
                alt.Tooltip("Destination", title="Recipient"),
                alt.Tooltip("AmountSafe", title="Amount ($)", format=",.2f"),
                alt.Tooltip("RiskSafe", title="Risk Score", format=".2%"),
                alt.Tooltip("Status", title="Verdict")
            ]
        ).properties(
            height=300,
            background="#0e1117"
        ).configure_view(strokeWidth=0)

        st.altair_chart(network_chart, use_container_width=True)
        st.caption("Figure 2: Transaction network — bubble size reflects transaction amount; colour encodes fraud probability.")
    else:
        st.info("📊 Analyse at least 2 transactions (or process a batch) to populate the network graph.")
    
    st.divider()
    
    # Feature Importance
    st.markdown("#### 2. Feature Importance Analysis")
    st.info("This chart explains **which factors** the model considers most important when detecting fraud.")
    
    try:
        current_dir = os.path.dirname(os.path.abspath(__file__))
        model_path = os.path.join(current_dir, '..', 'models', 'credit_fraud.pkl')
        model = joblib.load(model_path)
        
        # Map cryptic feature names to human-readable labels
        feature_map = {
            "step": "Time of Day (Hour)",
            "type": "Transaction Type",
            "amount": "Transaction Amount",
            "oldbalanceOrg": "Sender Initial Balance",
            "newbalanceOrig": "Sender Final Balance",
            "oldbalanceDest": "Receiver Initial Balance",
            "newbalanceDest": "Receiver Final Balance",
            "isFlaggedFraud": "System Flag (High Value)",
            "errorBalanceOrig": "Sender Balance Discrepancy",
            "errorBalanceDest": "Receiver Balance Discrepancy",
            "hour": "Hour of Day",
            "day_of_week": "Day of Week",
            "is_night": "Is Night Transaction",
            "amount_log": "Log(Amount)",
            "balance_dest_ratio": "Amount/Dest Balance Ratio"
        }
        
        # Handle dict format vs new format
        if isinstance(model, dict):
            # It's our new dict format
            feature_names = model.get('features', ["step", "type", "amount", "oldbalanceOrg", "newbalanceOrig", "oldbalanceDest", "newbalanceDest", "isFlaggedFraud"])
            
            # Check if pre-calculated importances exist (for VotingClassifier)
            if 'feature_importances' in model:
                importances = model['feature_importances']
            elif hasattr(model.get('model'), 'feature_importances_'):
                 importances = model['model'].feature_importances_
            else:
                 importances = np.zeros(len(feature_names))

        else:
             # Fallback for old list
            feature_names = ["step", "type", "amount", "oldbalanceOrg", "newbalanceOrig", "oldbalanceDest", "newbalanceDest", "isFlaggedFraud"]
            if hasattr(model, 'feature_importances_'):
                importances = model.feature_importances_
            else:
                importances = np.zeros(len(feature_names))

        # Ensure lengths match
        if len(feature_names) != len(importances):
             # Try resizing or create defaults if mismatch
             min_len = min(len(feature_names), len(importances))
             feature_names = feature_names[:min_len]
             importances = importances[:min_len]

        feat_df = pd.DataFrame({
            "Feature": [feature_map.get(f, f) for f in feature_names], 
            "Importance": importances
        }).sort_values(by="Importance", ascending=False).head(7)
            
        # Horizontal Bar Chart for better readability
        bars = alt.Chart(feat_df).mark_bar(color='#fca31180').encode(
            x=alt.X('Importance', title='Impact Power (0-1)'),
            y=alt.Y('Feature', sort='-x', title=''),
            tooltip=['Feature', alt.Tooltip('Importance', format='.1%')]
        )
        
        text = bars.mark_text(
            align='left',
            baseline='middle',
            dx=3,  # Nudges text to right so it doesn't overlap with the bar
            color='white'
        ).encode(
            text=alt.Text('Importance', format='.1%')
        )
        
        st.altair_chart((bars + text).properties(height=300), width="stretch")
        st.caption("Figure 1: Top factors driving the AI's fraud decisions.")

    except Exception as e:
        st.warning(f"Could not load feature importance: {e}")
        # Fallback static data if file not found (for demo purposes if running remotely without pkl)
        st.warning("Showing theoretical importance distribution (Model file unavailable).")
        data = {"oldbalanceOrg": 0.35, "amount": 0.25, "newbalanceOrig": 0.20, "type": 0.10, "step": 0.05, "oldbalanceDest": 0.03, "newbalanceDest": 0.02}
        st.bar_chart(data)

with tab1:
    # --- Quick Demo Buttons ---
    st.markdown("##### ⚡ Quick Demo")
    _TYPES_LIST = ["Cash In", "Cash Out", "Debit", "Payment", "Transfer"]
    demo_col1, demo_col2, demo_col3 = st.columns([1, 1, 4])
    with demo_col1:
        if st.button("🚨 Load Fraud Example", key="btn_fraud_demo"):
            # Realistic CASH_OUT fraud: sender drains account, receiver gets nothing logged
            st.session_state.demo_sender = "C6271823"
            st.session_state.demo_receiver = "C721638"
            st.session_state.demo_step = 206
            st.session_state.demo_amount = 245000.00
            st.session_state.demo_type_idx = 1   # Cash Out
            st.session_state.demo_oldbalanceorg = 245000.00
            st.session_state.demo_newbalanceorg = 0.00
            st.session_state.demo_oldbalancedest = 0.00
            st.session_state.demo_newbalancedest = 0.00
            st.toast("🚨 Fraud example loaded!", icon="🚨")
    with demo_col2:
        if st.button("✅ Load Safe Example", key="btn_safe_demo"):
            st.session_state.demo_sender = "C1239821"
            st.session_state.demo_receiver = "M928371"
            st.session_state.demo_step = 12
            st.session_state.demo_amount = 450.00
            st.session_state.demo_type_idx = 3   # Payment
            st.session_state.demo_oldbalanceorg = 8000.00
            st.session_state.demo_newbalanceorg = 7550.00
            st.session_state.demo_oldbalancedest = 300.00
            st.session_state.demo_newbalancedest = 750.00
            st.toast("✅ Safe example loaded!", icon="✅")

    st.markdown(
        "<div style='font-size:0.75rem;color:#6b7280;margin-bottom:8px'>Use the buttons above to pre-fill all fields with a known fraud or safe transaction for a quick demo.</div>",
        unsafe_allow_html=True
    )
    st.divider()

    with st.container():
        col_src, col_dest = st.columns(2)

        with col_src:
            st.subheader("📤 Source Account")
            sender_name = st.text_input("Account ID", value=st.session_state.demo_sender, placeholder="Ex: C12345678", key="input_sender")
            oldbalanceorg = st.number_input("Current Balance ($)", min_value=0.0, value=st.session_state.demo_oldbalanceorg, format="%.2f", key="input_oldbalanceorg")
            newbalanceorg = st.number_input("Post-Transaction Balance ($)", min_value=0.0, value=st.session_state.demo_newbalanceorg, format="%.2f", key="input_newbalanceorg")

        with col_dest:
            st.subheader("📥 Destination Account")
            receiver_name = st.text_input("Recipient ID", value=st.session_state.demo_receiver, placeholder="Ex: M12345678", key="input_receiver")
            oldbalancedest = st.number_input("Recipient Balance Before ($)", min_value=0.0, value=st.session_state.demo_oldbalancedest, format="%.2f", key="input_oldbalancedest")
            newbalancedest = st.number_input("Recipient Balance After ($)", min_value=0.0, value=st.session_state.demo_newbalancedest, format="%.2f", key="input_newbalancedest")

    st.divider()

    col_tx, col_meta = st.columns(2)

    with col_tx:
        st.subheader("💸 Transaction Details")
        amount = st.number_input("Transaction Amount ($)", min_value=0.0, value=st.session_state.demo_amount, format="%.2f", key="input_amount")
        types_map = {"Cash In": 0, "Cash Out": 1, "Debit": 2, "Payment": 3, "Transfer": 4}
        selected_type = st.selectbox("Transaction Type", list(types_map.keys()), index=st.session_state.demo_type_idx, key="input_type")
        types = types_map[selected_type]

    with col_meta:
        st.subheader("⚙️ Meta-Data")
        step = st.slider("Processing Time (Hours)", 0, 744, value=st.session_state.demo_step, key="input_step")
        isflaggedfraud = 1 if amount >= 200000 else 0
        st.caption(f"System Auto-Flag: {'🚩 High Value' if isflaggedfraud else '✅ Standard'}")
    
    # Analysis Button (Single)
    st.markdown("###")
    if st.button("🔍 ANALYZE TRANSACTION"):
        values = {
            "step": step,
            "types": types,
            "amount": amount,
            "oldbalanceorig": oldbalanceorg,
            "newbalanceorig": newbalanceorg,
            "oldbalancedest": oldbalancedest,
            "newbalancedest": newbalancedest,
            "isflaggedfraud": isflaggedfraud
        }
        
        try:
            # Animation & Progress Simulation
            with st.status("Initiating Protocol...", expanded=True) as status:
                st.write("🔌 Connecting to Secure Server...")
                time.sleep(0.6)
                st.write("🧠 Engaging Penta-Ensemble Model...")
                time.sleep(0.6)
                st.write("🔍 Scanning Transaction Patterns...")
                time.sleep(0.6)
                
                api_url = os.getenv("API_URL", "http://127.0.0.1:8000")
                res = re.post(f"{api_url}/predict", json=values)
                resp = res.json()
                status.update(label="Analysis Complete", state="complete", expanded=False)
                
            prediction = resp["result"]
            probability = resp["probability"]
            
            # Add to History
            st.session_state.history.append({
                "Timestamp": time.strftime("%H:%M:%S"),
                "Source": sender_name,
                "Destination": receiver_name,
                "Amount": amount,
                "Risk Score": probability,
                "Status": "Fraud" if prediction == "Fraudulent" else "Safe"
            })
            
            # Result Display
            st.markdown("---")
            st.subheader("📊 Examination Report")
            
            r_col1, r_col2, r_col3 = st.columns(3)
            
            with r_col1:
                st.metric("Risk Assessment", prediction, delta="CRITICAL" if prediction == "Fraudulent" else "SECURE", delta_color="inverse")
                
            with r_col2:
                st.metric("Fraud Probability", f"{probability:.2%}")
                st.progress(probability)
                
            with r_col3:
                confidence = max(probability, 1-probability)
                st.metric("Model Confidence", f"{confidence:.2%}")

            st.metric("Model Framework", "Heterogeneous Penta-Ensemble (RF, XGB, Ada, LR, MLP)")
            
            # detailed alert
            if prediction == "Fraudulent":
                st.error(f"""
                **🚨 SECURITY ALERT TRIGGERED**
                
                This transaction exhibits patterns consistent with financial fraud.
                - **Type:** {selected_type}
                - **Amount:** ${amount:,.2f}
                - **Risk Score:** {probability:.4f}
                
                **Action Required:** Immediate freeze recommended.
                """)
                st.toast("🚨 Threat Detected!", icon="🚨")
            else:
                st.success(f"""
                **✅ TRANSACTION CLEARED**
                
                No anomalies detected. Proceed with processing.
                - **Type:** {selected_type}
                - **Amount:** ${amount:,.2f}
                """)
                st.toast("✅ Verification Successful", icon="✅")
                
        except Exception as e:
            st.error(f"❌ Connection Error: Ensure Backend is running on Port 8000. Details: {e}")
        
        # ---------------------------
        # Explainable AI (SHAP)
        # ---------------------------
        if 'prediction' in locals(): # Only run if prediction succeeded
            st.markdown("---")
            st.subheader("🤖 AI Explanation (SHAP)")
            st.info("Why did the model make this decision? This chart shows how each feature pushed the risk score higher (Red) or lower (Blue).")
            
            try:
                # Call Explain Endpoint
                api_url = os.getenv("API_URL", "http://127.0.0.1:8000")
                exp_res = re.post(f"{api_url}/explain", json=values)
                if exp_res.status_code == 200:
                    explanation = exp_res.json()
                    
                    if "error" in explanation:
                        st.warning(f"Explanation unavailable: {explanation['error']}")
                    else:
                        shap_values = explanation["shap_values"]
                        feature_names_list = explanation["feature_names"]
                        
                        # Create DataFrame for Altair
                        # We map feature names to human-readable names using feature_map (defined earlier, need fallback here)
                        map_local = {
                            "step": "Time of Day (Hour)",
                            "type": "Transaction Type",
                            "amount": "Transaction Amount",
                            "oldbalanceOrg": "Sender Initial Balance",
                            "newbalanceOrig": "Sender Final Balance",
                            "oldbalanceDest": "Receiver Initial Balance",
                            "newbalanceDest": "Receiver Final Balance",
                            "isFlaggedFraud": "System Flag (High Value)",
                            "errorBalanceOrig": "Sender Balance Discrepancy",
                            "errorBalanceDest": "Receiver Balance Discrepancy",
                            "hour": "Hour of Day",
                            "day_of_week": "Day of Week",
                            "is_night": "Is Night Transaction",
                            "amount_log": "Log(Amount)",
                            "balance_dest_ratio": "Amount/Dest Balance Ratio"
                        }
                        readable_names = [map_local.get(f, f) for f in feature_names_list]
                        
                        explain_df = pd.DataFrame({
                            "Feature": readable_names,
                            "Impact": shap_values,
                            "Color": ["Increase Risk" if x > 0 else "Decrease Risk" for x in shap_values]
                        })
                        
                        # Filter for top impactful features (Top 8 absolute impact)
                        explain_df["AbsImpact"] = explain_df["Impact"].abs()
                        explain_df = explain_df.sort_values(by="AbsImpact", ascending=False).head(8)
                        
                        # Chart
                        shap_chart = alt.Chart(explain_df).mark_bar().encode(
                            x=alt.X('Impact', title='Impact on Fraud Probability'),
                            y=alt.Y('Feature', sort='-x', title=''),
                            color=alt.Color('Color', scale=alt.Scale(domain=['Increase Risk', 'Decrease Risk'], range=['#ff4b4b', '#1f77b4'])),
                            tooltip=['Feature', 'Impact']
                        ).properties(height=350)
                        
                        st.altair_chart(shap_chart, width="stretch")
                        
                else:
                    st.warning("⚠️ Could not retrieve explanation from server.")
            except Exception as e:
                st.warning(f"Could not generate explanation: {e}")

with tab2:
    st.subheader("📂 Batch Transaction Processor")
    st.verified = False
    uploaded_file = st.file_uploader("Upload CSV File (Requires: step, type, amount, etc.)", type=["csv"])
    
    # Template download
    sample_csv = pd.DataFrame({
        "step": [1, 1],
        "type": ["PAYMENT", "TRANSFER"],
        "amount": [100.0, 500000.0],
        "nameOrig": ["C12345", "C67890"],
        "oldbalanceOrg": [1000.0, 500000.0],
        "newbalanceOrig": [900.0, 0.0],
        "nameDest": ["M12345", "M67890"],
        "oldbalanceDest": [0.0, 0.0],
        "newbalanceDest": [0.0, 0.0]
    }).to_csv(index=False).encode('utf-8')
    
    st.download_button("download sample.csv", sample_csv, "sample_batch.csv", "text/csv")
    
    st.divider()

    if uploaded_file is not None:
        try:
            import pandas as pd
            df_upload = pd.read_csv(uploaded_file)
            st.dataframe(df_upload.head())
            
            col_opts, col_btn = st.columns([2, 1])
            with col_opts:
                include_explanations = st.checkbox("Include AI Explanations (Slower but detailed)", value=True)
            
            with col_btn:
                start_btn = st.button("🚀 Process Batch")

            if start_btn:
                progress_bar = st.progress(0)
                status_text = st.empty()
                
                results_list = []
                
                processed_count = 0
                fraud_count = 0
                update_interval = max(1, len(df_upload) // 20)
                
                for index, row in df_upload.iterrows():
                    # Feature Mapping
                    row_lower = {key.lower(): val for key, val in row.items()}
                    type_str = str(row_lower.get('type', 'Payment')).strip().title()
                    types_map = {"Cash In": 0, "Cash Out": 1, "Debit": 2, "Payment": 3, "Transfer": 4}
                    type_val = types_map.get(type_str, 3) 
                    
                    values = {
                        "step": int(row_lower.get('step', 1)),
                        "types": type_val,
                        "amount": float(row_lower.get('amount', 0.0)),
                        "oldbalanceorig": float(row_lower.get('oldbalanceorg', 0.0)),
                        "newbalanceorig": float(row_lower.get('newbalanceorig', 0.0)),
                        "oldbalancedest": float(row_lower.get('oldbalancedest', 0.0)),
                        "newbalancedest": float(row_lower.get('newbalancedest', 0.0)),
                        "isflaggedfraud": 1 if float(row_lower.get('amount', 0)) >= 200000 else 0
                    }
                    
                    top_factor = "N/A"
                    
                    try:
                        # 1. Predict
                        api_url = os.getenv("API_URL", "http://127.0.0.1:8000")
                        res = re.post(f"{api_url}/predict", json=values)
                        if res.status_code == 200:
                            resp = res.json()
                            pred = resp["result"]
                            prob = resp["probability"]
                            
                            # 2. Explain (if requested)
                            if include_explanations and (pred == "Fraudulent" or pred == "Not Fraudulent"):
                                try:
                                    exp_res = re.post(f"{api_url}/explain", json=values)
                                    if exp_res.status_code == 200:
                                        explanation = exp_res.json()
                                        if "error" not in explanation:
                                            shap_vals = explanation["shap_values"]
                                            feat_names = explanation["feature_names"]
                                            zipped = sorted(zip(feat_names, shap_vals), key=lambda x: x[1], reverse=True)
                                            map_local = {
                                                "amount": "Transaction Amount",
                                                "oldbalanceOrg": "Sender Initial Balance",
                                                "errorBalanceOrig": "Sender Balance Discrepancy",
                                                "errorBalanceDest": "Receiver Balance Discrepancy",
                                                "newbalanceOrig": "Sender Final Balance",
                                                "oldbalanceDest": "Receiver Initial Balance",
                                                "newbalanceDest": "Receiver Final Balance"
                                            }
                                            if zipped[0][1] > 0:
                                               friendly_name = map_local.get(zipped[0][0], zipped[0][0])
                                               top_factor = f"🚨 {friendly_name} (+{zipped[0][1]:.2f})"
                                            else:
                                               top_factor = "✅ Low Risk Profile"
                                except:
                                    pass

                            # Store Result
                            results_list.append({
                                "ID": str(row_lower.get('nameorig', f'TXN-{index}')),
                                "Amount": f"${values['amount']:,.2f}",
                                "Status": "Fraud" if pred == "Fraudulent" else "Safe",
                                "Risk Score": f"{prob:.2%}",
                                "Top Risk Factor": top_factor
                            })
                            
                            # Add to Session History if fraud (or limit size)
                            if pred == "Fraudulent" or len(st.session_state.history) < 100:
                                st.session_state.history.append({
                                    "Timestamp": time.strftime("%H:%M:%S"),
                                    "Source": str(row_lower.get('nameorig', 'Unknown')),
                                    "Destination": str(row_lower.get('namedest', 'Unknown')),
                                    "Amount": values['amount'],
                                    "Risk Score": prob,
                                    "Status": "Fraud" if pred == "Fraudulent" else "Safe"
                                })
                            
                            if pred == "Fraudulent":
                                fraud_count += 1
                        
                    except Exception as e:
                        print(f"Error processing row {index}: {e}")
                        
                    processed_count += 1
                    
                    if processed_count % update_interval == 0 or processed_count == len(df_upload):
                        progress_bar.progress(processed_count / len(df_upload))
                        status_text.text(f"Processing {processed_count}/{len(df_upload)} transactions...")
                    
                st.success(f"✅ Batch Complete! {len(df_upload)} scanned, {fraud_count} potential threats found.")
                
                # Show Results Code
                res_df = pd.DataFrame(results_list)
                st.dataframe(res_df.style.apply(lambda x: ['background-color: #ff4b4b' if x.Status == 'Fraud' else '' for _ in x], axis=1), width="stretch")
                
                # Download Result
                csv_res = res_df.to_csv(index=False).encode('utf-8')
                st.download_button(
                    "📥 Download Analysis Report",
                    csv_res,
                    "fraud_analysis_report.csv",
                    "text/csv",
                    key='download-batch-res'
                )
                
        except Exception as e:
            st.error(f"Error reading file: {e}")

# Session History Section
if st.session_state.history:
    st.divider()
    st.markdown("### 📜 Session Activity Log")
    
    # History Controls
    col_hist_1, col_hist_2 = st.columns([1, 6])
    with col_hist_1:
        if st.button("🗑️ Clear History"):
            st.session_state.history = []
            st.rerun()
    with col_hist_2:
        if st.button("↩️ Undo Last"):
            if len(st.session_state.history) > 0:
                st.session_state.history.pop()
                st.rerun()
    
    # Convert to DataFrame for display
    import pandas as pd
    df = pd.DataFrame(st.session_state.history)
    
    # Style the dataframe (Highlight Fraud rows)
    def highlight_fraud(row):
        return ['background-color: #ff4b4b' if row.Status == 'Fraud' else '' for _ in row]
    
    st.dataframe(df.style.apply(highlight_fraud, axis=1), width="stretch")
    
    # Advanced Altair Chart
    
    st.markdown("### 📈 Risk Trend Analysis")
    
    # Create valid dataframe for plotting
    chart_df = df.copy()
    chart_df["Order"] = range(len(chart_df))
    
    # Base Chart
    base = alt.Chart(chart_df).encode(
        x=alt.X('Order', axis=alt.Axis(title='Transaction Sequence')),
        tooltip=['Timestamp', 'Source', 'Amount', 'Risk Score', 'Status']
    )

    # Line Layer
    line = base.mark_line(color='#fca311', strokeWidth=3).encode(
        y=alt.Y('Risk Score', scale=alt.Scale(domain=[0, 1]), axis=alt.Axis(title='Fraud Probability'))
    )
    
    # Area Layer (Gradient)
    area = base.mark_area(opacity=0.3, color='#fca311').encode(
        y=alt.Y('Risk Score')
    )
    
    # Points Layer (for hover interaction)
    points = base.mark_circle(size=60, color='white').encode(
        y=alt.Y('Risk Score')
    )
    
    # Threshold Line (Red line at 0.5)
    threshold = alt.Chart(pd.DataFrame({'y': [0.5]})).mark_rule(color='red', strokeDash=[5, 5]).encode(
        y='y'
    )
    
    final_chart = (area + line + points + threshold).properties(
        height=300
    ).interactive()
    
    st.altair_chart(final_chart, width="stretch")
    
    # Download Button
    csv = df.to_csv(index=False).encode('utf-8')
    st.download_button(
        "📥 Download Activity Log",
        csv,
        "fraud_guard_log.csv",
        "text/csv",
        key='download-csv'
    )

\end{lstlisting}

